\documentclass[11pt]{article}

\usepackage[utf8]{inputenc}
\usepackage[T1]{fontenc}
\usepackage{geometry}
\usepackage{amsmath}
\usepackage{amssymb}
\usepackage{booktabs}
\usepackage{hyperref}
\usepackage{xurl}
\usepackage{xcolor}
\usepackage{listings}
\usepackage{enumitem}
\usepackage{parskip}
\usepackage{graphicx}
\graphicspath{{../images/}}

\geometry{margin=1in}

\lstset{
  language=Java,
  basicstyle=\ttfamily\small,
  keywordstyle=\color{blue},
  commentstyle=\color{gray},
  stringstyle=\color{teal},
  showstringspaces=false,
  columns=fullflexible,
  frame=single,
  framerule=0.2pt,
  breaklines=true
}

\setlist[itemize]{topsep=0.3em,itemsep=0.2em}
\setlist[enumerate]{topsep=0.3em,itemsep=0.2em}

% Simple per-section source line.
\newcommand{\notesources}[1]{\par\noindent{\small\color{gray}\textit{Sources:} \sloppy #1\par}}

% Unit-wise source strings for this subject (used by slides/notes).
% Path is relative to the lecture `latex/` folder (the compilation CWD).
% OOPS with Java (CSEG1044) — unit-wise sources
%
% These are short, student-facing reference strings intended to be used with:
%   - \slidesources{...}  (in beamer slides)
%   - \notesources{...}   (in lecture notes)
%
% Style inspiration: other subjects in My-Subjects/ use semicolon-separated sources
% and include an "accessed" date for web documentation.

\newcommand{\OOPSUnitISources}{Schildt, \textit{Java: A Beginner's Guide} (10e), McGraw-Hill (Java fundamentals + OOP basics).; Oracle Java Tutorials: Getting Started + Language Basics + Classes and Objects (accessed \today).; Java SE API docs (e.g., \texttt{String}, \texttt{Scanner}) (accessed \today).}

\newcommand{\OOPSUnitIISources}{Schildt, \textit{Java: The Complete Reference} (12e), McGraw-Hill (classes, inheritance, packages, interfaces).; Bloch, \textit{Effective Java} (3e), Addison-Wesley, 2018 (design choices; interfaces vs abstract classes).; Oracle Java Tutorials: Classes and Objects + Interfaces + Packages (accessed \today).; Java SE API docs (e.g., \texttt{Comparable}, \texttt{Runnable}) (accessed \today).}

\newcommand{\OOPSUnitIIISources}{Schildt, \textit{Java: The Complete Reference} (12e) (nested/inner classes, exceptions, threads, I/O).; Oracle Java Tutorials: Nested Classes + Exceptions + Concurrency + Basic I/O (accessed \today).; Java SE API docs (e.g., \texttt{Thread}, \texttt{AutoCloseable}) (accessed \today).}

\newcommand{\OOPSUnitIVSources}{Schildt, \textit{Java: The Complete Reference} (12e) (generics, lambdas, Swing, JDBC).; Bloch, \textit{Effective Java} (3e) (generics + lambdas).; Oracle Java Tutorials: Generics + Lambda Expressions + Swing + JDBC Basics (accessed \today).; Java SE API docs (e.g., \texttt{java.util.function}, \texttt{javax.swing}, \texttt{java.sql}) (accessed \today).}

\newcommand{\OOPSUnitVSources}{Schildt, \textit{Java: The Complete Reference} (12e) (Collections Framework + wrapper classes).; Oracle Java docs: Collections Framework overview (accessed \today).; Java SE API docs (\texttt{java.util} collections; wrapper classes in \texttt{java.lang}) (accessed \today).}

\newcommand{\OOPSUnitVISources}{Git SCM: \textit{Pro Git} (2e) (version control workflow), accessed \today.; GitHub Docs (repositories, issues, pull requests), accessed \today.; Oracle Java Tutorials: Swing + JDBC (accessed \today).; JUnit 5 User Guide (accessed \today).; Apache Maven Project: Getting Started (accessed \today).; Gradle User Manual: Getting Started (accessed \today).; UML class diagrams: UML 2.x overview/spec (accessed \today).}

\newcommand{\OOPSMiscWhyInterfacesSources}{Schildt, \textit{Java: The Complete Reference} (12e) (interfaces).; Bloch, \textit{Effective Java} (3e) (prefer interfaces where appropriate).; Oracle Java Tutorials: Interfaces (accessed \today).; Java SE API docs (e.g., \texttt{Comparable}, \texttt{Runnable}) (accessed \today).}



\title{Object Oriented Programming Lab (CSEG1044)\\Experiment-13\&14: File I/O + Generics (Student Handout --- Before Submission)}
\author{Tofik Ali}
\date{\today}

\begin{document}
\maketitle
\tableofcontents

\section{About this handout}
This handout specifies the required programs for file handling and generics.
Write and test your own implementation.

\section{Experiment 13 (Syllabus)}
\begin{itemize}[leftmargin=*]
  \item \textbf{Objective:} To perform file operations using Java I/O streams.
  \item \textbf{Problem Statement:} Develop programs using \texttt{FileReader} and \texttt{FileWriter}.
  \item \textbf{Expected Outcome:} Students will handle persistent data.
\end{itemize}

\section{Experiment 14 (Syllabus)}
\begin{itemize}[leftmargin=*]
  \item \textbf{Objective:} To implement generic classes and methods.
  \item \textbf{Problem Statement:} Design programs using generics.
  \item \textbf{Expected Outcome:} Students will understand type safety.
\end{itemize}

\section{What you must do (minimum requirements)}

\subsection{Task A --- File I/O (Exp-13)}
Create \texttt{FileIOStreamsDemo.java}:
\begin{enumerate}[leftmargin=*]
  \item Write 3 lines to \texttt{input.txt} using \texttt{FileWriter}.
  \item Read and print the file contents using \texttt{FileReader}.
  \item Copy \texttt{input.txt} to \texttt{copy.txt} using a character buffer.
  \item Print how many characters were copied.
  \item Use try-with-resources and handle \texttt{IOException}.
\end{enumerate}

\subsection{Task B --- Generics basics (Exp-14)}
Create \texttt{GenericsDemo.java}:
\begin{enumerate}[leftmargin=*]
  \item Create a generic class \texttt{Box<T>} that stores a single value and returns it via \texttt{get()}.
  \item Create a generic method \texttt{printTwice(T x)} and call it with different types.
  \item Demonstrate at least:
  \begin{itemize}[leftmargin=*]
    \item \texttt{Box<Integer>}
    \item \texttt{Box<String>}
  \end{itemize}
  \item Optional: create a generic \texttt{Pair<K,V>} and print it.
\end{enumerate}

\section{Hints (read only if stuck)}
\begin{itemize}[leftmargin=*]
  \item Use try-with-resources to auto-close \texttt{Reader/Writer}.
  \item Relative file paths are resolved from your program's working directory (often your project folder).
  \item Generics prevent many \texttt{ClassCastException} issues by enforcing type checks at compile time.
  \item Avoid raw types (use \texttt{Box<String>} instead of \texttt{Box}).
\end{itemize}

\section{What to submit (MANDATORY)}
Submit a single zipped folder containing:
\begin{itemize}[leftmargin=*]
  \item \texttt{FileIOStreamsDemo.java}
  \item \texttt{GenericsDemo.java}
  \item sample output for both programs
  \item observations (5--10 lines): what try-with-resources does; where generics helped
\end{itemize}

\notesources{\OOPSUnitIIISources\ \OOPSUnitIVSources}
\end{document}
